\section{Pengenalan Server Web (Apache/Nginx) dan Peran PHP}

\textbf{Server web} bertugas menerima HTTP request dari browser dan mengembalikan response (file HTML, CSS, JS, gambar, atau output PHP). Dua server web paling populer: \textbf{Apache HTTP Server} (menggunakan \code{.htaccess} untuk konfigurasi per-direktori; modul \code{mod\_php} menjalankan PHP langsung) dan \textbf{Nginx} (event-driven, performa tinggi untuk koneksi bersamaan; menggunakan PHP-FPM sebagai proses terpisah). Apache lebih mudah dikonfigurasi untuk pemula; Nginx lebih efisien untuk traffic tinggi dan sering digunakan sebagai reverse proxy \cite{mdnDeploy}.

\begin{table}[ht]
\centering
\begin{tabular}{|P{3cm}|P{5cm}|P{5cm}|}
\hline
\textbf{Aspek} & \textbf{Apache} & \textbf{Nginx} \\
\hline
Arsitektur & Process/thread per request & Event-driven, async \\
\hline
PHP handler & mod\_php (built-in) & PHP-FPM (terpisah) \\
\hline
Konfigurasi & .htaccess (per-dir) & conf file (terpusat) \\
\hline
Performa & Baik untuk traffic sedang & Sangat baik untuk traffic tinggi \\
\hline
Reverse proxy & Bisa (mod\_proxy) & Built-in, sangat populer \\
\hline
Rewrite URL & mod\_rewrite & \code{try\_files} directive \\
\hline
\end{tabular}
\caption{Perbandingan Apache dan Nginx}
\end{table}

